\documentclass[titlepage,a4paper]{jsarticle}
\usepackage{../../../sty/import}% 各種パッケージインポート
\usepackage{../../../sty/title}% タイトルページの変更

%% タイトルページの変数
% レポートタイトル
\title{原子力材料と核燃料}
% 提出日
\expdate{\today}
% 科目名
\subject{原子力材料と核燃料}
% 分野
\class{情報経営システム工学分野}
% 学年
\grade{M1}
% 学籍番号
\mynumber{24336488}
% 記述者
\author{本間三暉}

\begin{document}
% titleページ作成
\maketitle
\section{鉄鋼組織と冷却速度の関係を示せ}
鉄鋼をオーステナイト（$\gamma$相）の状態から冷却すると、冷却速度によって
生成する組織が異なる。これは、冷却速度によって炭素原子が拡散するための
時間が変化するためである。

\subsection{徐冷した場合}

オーステナイトをゆっくりと冷却すると、炭素原子が十分に拡散できるため、
パーライトが形成される。パーライトは、フェライト（$\alpha$相）と
セメンタイト（Fe$_3$C）が層状に並んだ組織である。また、鋼の炭素濃度が
共析組成と異なる場合には、初析フェライトまたは初析セメンタイトが加わる。

パーライトの層間隔は冷却速度によって変化する。冷却速度が特に遅い場合には
粗大パーライトとなり、比較的速く冷却した場合には微細パーライトとなる。
微細パーライトは粗大パーライトよりも組織が細かいため、強度が高くなる。

\subsection{中程度の速度で冷却した場合}

オーステナイトを中程度の速度で冷却すると、ベイナイトが形成される。
ベイナイトは、フェライトと板状または針状のセメンタイトから構成される
組織である。パーライトよりも微細な組織を持つため、一般にパーライトより
強度が高い。一方で、マルテンサイトと比較すると強度は低いが、延性は高い。

\subsection{急冷した場合}

オーステナイトを急冷すると、炭素原子が拡散するための時間がなくなる。
そのため、拡散を伴わない変態によってマルテンサイトが形成される。
マルテンサイトは、炭素を過飽和に含む体心正方晶
（Body-Centered Tetragonal: BCT）構造を持つ。

マルテンサイトは非常に高い強度と硬さを持つ一方で、延性が低く、
脆くなりやすい。したがって、実際の材料では、急冷後に再加熱する
焼戻し処理を行うことがある。焼戻しによって、フェライト中に非常に微細な
セメンタイト粒子が分散した焼戻しマルテンサイトが形成される。
焼戻しを行うことで、マルテンサイトの高い強度をある程度維持しながら、
延性や靱性を改善できる。

\subsection{強度および延性の変化}

各組織の強度は、一般に次の順に高くなる。

\begin{equation}
  \text{粗大パーライト}
  <
  \text{微細パーライト}
  <
  \text{ベイナイト}
  <
  \text{焼戻しマルテンサイト}
  <
  \text{マルテンサイト}
\end{equation}

一方、延性はこれと反対の順序となり、マルテンサイトが最も低く、
粗大パーライトが最も高い傾向を示す。したがって、冷却速度を速くすると
生成する組織は微細化し、強度および硬さは増加するが、延性は低下する。
ただし、焼戻しマルテンサイトは急冷によって直接生成する組織ではなく、
急冷によって得られたマルテンサイトを再加熱することで形成される。

以上より、鉄鋼材料では、冷却速度および急冷後の熱処理条件を調整することで、
強度、硬さ、延性および靱性のバランスを制御できる。
% 参考文献
\begin{thebibliography}{99}
\bibitem{} 原子力材料と核燃料 講義資料
\end{thebibliography}

\end{document}